% !TEX root = ../../../main.tex
% 来源：中学数学实验教材·第一册（代数）
% 章节：因式分解与余式定理 / 因式分解
% 原始文件：5.tex，行 633-664
% 类型：tikz
% ---
        \begin{tikzpicture}
     \begin{scope}
    \node (A) at (-.6,1.2) {$2$};
    \node (B) at (.6,1.2) {$-4$};
    \node (C) at (-.6,0) {$3$};
    \node (D) at (.6,0) {$3$};
    \draw[very thick] (A)--(D);
    \draw[very thick] (B)--(C);
    \node at (0,-.5){$-12+6=-6$};
    \node at (0,-1){（不合适）};
     \end{scope}   
     \begin{scope}[xshift=3cm]
        \node (A) at (-.6,1.2) {$2$};
        \node (B) at (.6,1.2) {$3$};
        \node (C) at (-.6,0) {$3$};
        \node (D) at (.6,0) {$-4$};
        \draw[very thick] (A)--(D);
        \draw[very thick] (B)--(C);
        \node at (0,-.5){$9-8=1$};
        \node at (0,-1){（不合适）};
         \end{scope} 
         \begin{scope}[xshift=6cm]
            \node (A) at (-.6,1.2) {$2$};
            \node (B) at (.6,1.2) {$-3$};
            \node (C) at (-.6,0) {$3$};
            \node (D) at (.6,0) {$4$};
            \draw[very thick] (A)--(D);
            \draw[very thick] (B)--(C);
            \node at (0,-.5){$-9+8=-1$};
            \node at (0,-1){（合适）};
             \end{scope} 
    \end{tikzpicture}
